\begin{abstract}
Unlabeled test data can improve the zero-shot classifier of a vision-language model, but realistic evaluations show that transductive and test-time adaptation methods often fall below the zero-shot baseline when only a few vocabulary classes are present in a batch, when streams are correlated, or when the backbone changes. We trace a common cause to how adaptation is parameterised: every class owns independent parameters estimated from its own test samples, so absent or rare classes are either left at a biased text anchor or fitted to samples of other classes. We propose \method (Graph-Regularized Anchor-displacement Field Transduction), which models the adapted prototypes as the text anchors plus a displacement field that is partially pooled over a semantic class graph built from the text encoder. A Gaussian Markov random field prior lets rare and absent classes borrow the displacement of their well-observed semantic neighbours, and an explicit class-presence model turns the vocabulary into a label-shift estimate instead of a fixed set of competitors. Inference is closed form and gradient free on cached features; the anchor-only shrinkage of prior work and a single global domain shift are the two extreme settings of one prior precision, whose strength can be chosen from the unlabeled batch by empirical Bayes. We show that absent classes receive the harmonic extension of the present-class displacements and that the posterior variance of every prototype is never larger than under independent shrinkage. The method is implemented with unit-tested closed-form updates and is evaluated under the realistic batch and online protocols on eleven benchmarks and ImageNet distribution shifts with several CLIP-style backbones; benchmark numbers are reported as the pre-registered experiments complete and no result in this draft is estimated.
\end{abstract}
